\documentclass[12pt,a4paper]{ctexart}

% 页面与基础
\usepackage[margin=2.5cm]{geometry}
\usepackage{setspace}
\usepackage{indentfirst}   % 让章节后首段也缩进
\setlength{\parindent}{2em}

\setstretch{1.2}

% 数学
\usepackage{amsmath,amssymb,bm}

% 图表
\usepackage{graphicx}
\usepackage{subcaption}
\usepackage{booktabs}
\usepackage{array}
\usepackage{float}
\usepackage{caption}

% 代码与链接
\usepackage{listings}
\usepackage{xcolor}
\usepackage{hyperref}
\hypersetup{
  colorlinks=true,
  linkcolor=black,
  urlcolor=blue,
  citecolor=blue
}

% 代码样式
\lstset{
  basicstyle=\ttfamily\small,
  keywordstyle=\color{blue},
  commentstyle=\color{teal},
  stringstyle=\color{orange},
  numbers=left,
  numberstyle=\tiny,
  stepnumber=1,
  frame=single,
  breaklines=true,
  tabsize=2
}



\begin{document}
% ===== 封面页 =====
\begin{titlepage}
  \centering
  \vspace*{0.16\textheight}
  {\zihao{2}\bfseries 物理信息神经网络（PINNs）实验报告\par}

  \vspace*{0.48\textheight}

  {\zihao{4}
  \begin{tabular}{rrc}
    姓名&： & 张春啸 \\
    学号&： & 2021013475 \\
    日期&： & 2026/3/30 \\
  \end{tabular}
  \par}

  \vfill
\end{titlepage}

% ===== 目录页 =====
\clearpage
\tableofcontents
\thispagestyle{plain}
\clearpage

\section{FDM数值方法}
为获得 Burgers 方程的高可信参考解，采用有限差分法（FDM）进行数值求解。具体做法是：在空间区间 $x\in[-1,1]$ 与时间区间 $t\in[0,1]$ 上构造均匀网格；对扩散项 $u_{xx}$ 使用二阶中心差分离散，对对流项 $u u_x$ 使用随速度符号切换的一阶迎风差分离散；时间推进采用显式欧拉格式，并在每个时间步强制施加边界条件 $u(-1,t)=u(1,t)=0$。初值取 $u(x,0)=-\sin(\pi x)$。最终得到的时空离散解作为 PINNs 训练后评估时的参考真值，用于计算相对 $L_2$ 误差与 MSE 等指标。

\begin{figure}[H]
  \centering
  \includegraphics[width=0.75\textwidth]{topic2_pinns/outputs_run_1e-2/base/ref_field.png}
  \caption{$\nu =0.01/\pi$ 下的FDM参考场}
\end{figure}

\section{标准PINNs方法}
标准 PINNs 的核心思想是：用神经网络 $u_\theta(x,t)$ 近似 PDE 的解函数，并通过自动微分将物理方程残差直接加入训练目标。本文采用多层全连接网络（MLP）作为近似器，输入为 $(x,t)$，输出为标量场值 $u$。\par

针对 Burgers 方程\par
\[
u_t + u u_x = \nu u_{xx},
\]\par
构造三部分损失：\par
\[
\mathcal{L}=\mathcal{L}_{ic}+\mathcal{L}_{bc}+\mathcal{L}_{pde},
\]\par
其中\par
\[
\mathcal{L}_{ic}=\frac{1}{N_{ic}}\sum_{i=1}^{N_{ic}}\left|u_\theta(x_i,0)-u_0(x_i)\right|^2,
\]\par
\[
\mathcal{L}_{bc}=\frac{1}{N_{bc}}\sum_{j=1}^{N_{bc}}\left(\left|u_\theta(-1,t_j)-0\right|^2+\left|u_\theta(1,t_j)-0\right|^2\right),
\]\par
\[
\mathcal{L}_{pde}=\frac{1}{N_f}\sum_{k=1}^{N_f}\left|u_t+u u_x-\nu u_{xx}\right|_{(x_k,t_k)}^2.
\]\par
这里 $u_t,u_x,u_{xx}$ 均由自动微分计算。训练时使用 Adam 优化器最小化总损失，并在采样点上不断迭代更新网络参数。最终在规则网格上评估预测解，与 FDM 参考解对比，计算相对 $L_2$ 误差和 MSE 作为性能指标。\par

控制方程：\par
\[
u_t + u u_x = \nu u_{xx}, \quad x\in[-1,1],\ t\in[0,1].
\]\par

\section{改进方法}
针对标准 PINNs 在小粘性系数下容易出现的收敛困难与误差偏大问题，本文采用了两种改进策略：自适应激活函数与残差自适应重采样（RAR）。\par

\subsection{自适应激活函数（ada）}
标准 MLP 常用固定激活 $\tanh(z)$。在本方法中，将其改为\par
\[
\tanh(\alpha z),
\]\par
其中 $\alpha$ 为可训练参数，并与网络权重共同优化。该设计能够动态调节网络非线性强度，使模型在训练过程中更灵活地适配不同尺度的解结构，尤其有助于提升对陡峭梯度区域的表达能力。

\subsection{残差自适应重采样（rar）}
标准 PINNs 常采用均匀随机采样 PDE 配点，可能导致高残差区域关注不足。RAR 的思路是周期性执行以下步骤：
\begin{enumerate}
  \item 在候选点池中随机采样；
  \item 计算各点 PDE 残差幅值；
  \item 选取残差最大的若干点（hard points）；
  \item 用这些点替换部分训练配点，强化后续训练。
\end{enumerate}\par
通过“哪里误差大就重点学哪里”，RAR 可以在一定程度上缓解采样与难点区域错配问题，提升训练稳定性。\par

本实验将以上两种方法分别与基线方法对比，评估它们在不同粘性系数下对相对 $L_2$ 误差和 MSE 的影响，并分析其有效性与局限性。\par

\section{实验结果与分析}

\subsection{可视化结果对比}

为直观比较模型在不同粘性系数下的表现，本文先给出两组参数下的解场对比图。每组包含 4 幅图：FDM 参考解、base、ada、rar。

\subsubsection{$\nu=0.01/\pi$}
\begin{figure}[H]
  \centering
  \begin{subfigure}{0.48\textwidth}
    \centering
    \includegraphics[width=\textwidth]{topic2_pinns/outputs_run_1e-2/base/ref_field.png}
    \caption{FDM}
  \end{subfigure}
  \hfill
  \begin{subfigure}{0.48\textwidth}
    \centering
    \includegraphics[width=\textwidth]{topic2_pinns/outputs_run_1e-2/base/pred_field.png}
    \caption{base}
  \end{subfigure}

  \vspace{0.6em}

  \begin{subfigure}{0.48\textwidth}
    \centering
    \includegraphics[width=\textwidth]{topic2_pinns/outputs_run_1e-2/ada/pred_field.png}
    \caption{ada}
  \end{subfigure}
  \hfill
  \begin{subfigure}{0.48\textwidth}
    \centering
    \includegraphics[width=\textwidth]{topic2_pinns/outputs_run_1e-2/rar/pred_field.png}
    \caption{rar}
  \end{subfigure}
  \caption{$\nu=0.01/\pi$ 下三种 PINNs 与 FDM 对比}
\end{figure}

\begin{figure}[H]
  \centering
  \begin{subfigure}{0.32\textwidth}
    \centering
    \includegraphics[width=\textwidth]{topic2_pinns/outputs_run_1e-2/base/error_field.png}
    \caption{base}
  \end{subfigure}
  \hfill
  \begin{subfigure}{0.32\textwidth}
    \centering
    \includegraphics[width=\textwidth]{topic2_pinns/outputs_run_1e-2/ada/error_field.png}
    \caption{ada}
  \end{subfigure}
  \hfill
  \begin{subfigure}{0.32\textwidth}
    \centering
    \includegraphics[width=\textwidth]{topic2_pinns/outputs_run_1e-2/rar/error_field.png}
    \caption{rar}
  \end{subfigure}
  \caption{$\nu=0.01/\pi$ 下三种方法的误差场对比}
\end{figure}

\subsubsection{$\nu=0.001/\pi$}
\begin{figure}[H]
  \centering
  \begin{subfigure}{0.48\textwidth}
    \centering
    \includegraphics[width=\textwidth]{topic2_pinns/outputs_run_1e-3/base/ref_field.png}
    \caption{FDM}
  \end{subfigure}
  \hfill
  \begin{subfigure}{0.48\textwidth}
    \centering
    \includegraphics[width=\textwidth]{topic2_pinns/outputs_run_1e-3/base/pred_field.png}
    \caption{base}
  \end{subfigure}

  \vspace{0.6em}

  \begin{subfigure}{0.48\textwidth}
    \centering
    \includegraphics[width=\textwidth]{topic2_pinns/outputs_run_1e-3/ada/pred_field.png}
    \caption{ada}
  \end{subfigure}
  \hfill
  \begin{subfigure}{0.48\textwidth}
    \centering
    \includegraphics[width=\textwidth]{topic2_pinns/outputs_run_1e-3/rar/pred_field.png}
    \caption{rar}
  \end{subfigure}
  \caption{$\nu=0.001/\pi$ 下三种 PINNs 与 FDM 对比}
\end{figure}

\begin{figure}[H]
  \centering
  \begin{subfigure}{0.32\textwidth}
    \centering
    \includegraphics[width=\textwidth]{topic2_pinns/outputs_run_1e-3/base/error_field.png}
    \caption{base}
  \end{subfigure}
  \hfill
  \begin{subfigure}{0.32\textwidth}
    \centering
    \includegraphics[width=\textwidth]{topic2_pinns/outputs_run_1e-3/ada/error_field.png}
    \caption{ada}
  \end{subfigure}
  \hfill
  \begin{subfigure}{0.32\textwidth}
    \centering
    \includegraphics[width=\textwidth]{topic2_pinns/outputs_run_1e-3/rar/error_field.png}
    \caption{rar}
  \end{subfigure}
  \caption{$\nu=0.001/\pi$ 下三种方法的误差场对比}
\end{figure}





\subsection{定量指标分析}

对应误差指标如下（rel\_l2 越小越好）：

\begin{table}[H]
  \centering
  \caption{两组粘性系数下三种方法误差对比}
  \begin{tabular}{c c c c}
    \toprule
    $\nu_{\text{input}}$ & 方法 & rel\_l2 & mse \\
    \midrule
    0.01  & base & 6.94496e-2 & 1.79318e-3 \\
    0.01  & ada  & \textbf{2.75934e-2} & \textbf{2.83071e-4} \\
    0.01  & rar  & 1.63941e-1 & 9.99209e-3 \\
    0.001 & base & \textbf{4.47943e-1} & \textbf{7.62785e-2} \\
    0.001 & ada  & 5.82598e-1 & 1.29031e-1 \\
    0.001 & rar  & 4.48329e-1 & 7.64101e-2 \\
    \bottomrule
  \end{tabular}
\end{table}

\subsection{结果讨论}

在 $\nu=0.01/\pi$ 条件下，ada 的结果最接近 FDM，误差明显低于 base 与 rar，说明自适应激活在该难度下能够提升拟合能力。  
当粘性系数减小到 $\nu=0.001/\pi$ 时，三种方法误差均显著上升，表现出强对流场景下 PINNs 的训练退化现象；其中 base 与 rar 接近，ada 退化更明显。  
总体来看，改进方法并非在所有参数区间都稳定优于基线，方法效果与问题刚性程度和训练超参数存在显著耦合，需要针对小粘性场景进一步调参或引入更强的训练策略。


\end{document}
